
 \Cref{sync_maximal:thm:structure} and \cref{sync_max:cor:structure_letters}
 gave structural results linking the class of sync-maximal automata
 to completely reachable and completely distinguishable automata.
 Later, in \cref{sync_max:sec:SI}, we will give more precise
 structural results for classes of automata.
 For example, an automaton whose transformation semigroup contains
 a transitive group of units and where every letter has defect at most one is synchronizing
 if and only if it is sync-maximal, see \cref{sync_max:sec:defectone}.
 This statement applies to the the \Cerny automata from \cref{intro:fig:cerny}.
 
We note that deciding sync-maximality is \PSPACE-complete, see
the discussion in \cref{sync_max:decision:discussion}.

\IN{reset complexity}
By \cref{prelim:thm:sync_words_is_ideal}, languages of the form $\Syn(\mathcal A)$ are precisely the regular ideal languages.
Thus, shortest words in regular ideal languages corresponds
to shortest paths in accepting automata for sets of synchronizing words.
Now, investigating the length of the shortest words in regular ideal languages as a language-theoretic approach to the \v{C}ern\'{y} conjecture has been pursued and extended by several authors~\cite{DBLP:journals/ijfcs/AlmeidaR16,DBLP:journals/jalc/FrigeriR25,DBLP:journals/jalc/FrigeriR19,DBLP:journals/fuin/MaslennikovaR18,DBLP:journals/iandc/PribavkinaR11,ReisR16,DBLP:journals/aam/Rodaro18}, based on a method introduced by Rodero et al.~\cite{DBLP:journals/fuin/MaslennikovaR18,ReisR16,DBLP:journals/aam/Rodaro18} that even associates a strongly connected synchronizing automaton with each regular ideal language.
Maslennikova~\cite{DBLP:journals/corr/Maslennikova14,DBLP:journals/ijfcs/Maslennikova19} 
introduced the \emph{reset complexity} 
of a regular ideal language $I \subseteq \Sigma^*$ 
as the number of states of a smallest automaton~$\mathcal A$ 
such that $I = \Syn(\mathcal A)$.
For a sync-maximal automaton $\mathcal A$,
the gap between the size of an accepting automaton for $\Syn(\mathcal A)$, i.e., the state complexity of this set,
and the size of an automaton having $\Syn(\mathcal A)$ as its set of synchronizing words, i.e., its reset complexity,
is maximal possible (note that in this case, no automaton with fewer states than~$\mathcal A$ can have $\Syn(\mathcal A)$
as its set of synchronizing words).
In this sense, for sync-maximal automata the complexity of the set of synchronizing words is maximal.